\scanpage{7}

\[\rho(C,\delta)=\rho_3\rho_2\rho_1(C,\delta)\]

On définit de façon duale une \deleted{opération}\inserted{construction} $\sigma=\sigma_3\sigma_2\sigma_1$, assurant successivement $b'')$, $\deleted{b}a''_1$, $a''_2$.

On pose $K(C,\delta)=\sigma\rho(C,\delta)=\sigma_3\sigma_2\sigma_1\rho_3\rho_2\rho_1(C,\delta)$. Il a les propriétés a) à d) voulues.

\uline{NB} L'ordre dans lequel on a composé les \inserted{6 opérations} $\sigma_i$, \deleted{et les} $\rho_i$ n'a pas d'importance. On pourrait prendre p. ex
\[
K'(C,\delta)=\sigma_1\rho_3\sigma_2\rho_1\rho_2\sigma_3(C,\delta)\qquad !
\]
La chose ayant une signification intrinsèque est la limite des \deleted{\unsure{l'}} itérés transfinis de $K$....

\begin{center}\rule{1.1cm}{0.4pt}\end{center}

\uline{Conditions d'exactitude.}

On a donné \inserted{\unsure{trois paquets|tripliquet|triplet} de $\Delta$-foncteurs entre} des petites catégories dans $\Delta=\Delta_d=\Delta_{d',d''}$, \inserted{(petits)}
\[
f^{(i)}_\lambda:R^{\unsure{(i)}}_\lambda\unsure{-|\to}R^{\prime(i)}_\lambda\quad\inserted{(i=0,1,2)},\quad\text{on pose } f^i=(f_\lambda),\quad f=(f^{(0)},f^{(1)},f^{(2)})
\]

$\Delta_f$ est la sous-$2$-catégorie \inserted{$2$-}pleine de $\Delta$ formée des $E\in\mathrm{Ob}\,\Delta$ tels que, \deleted{le foncteur $\varphi\longmapsto\varphi\vert f$} pour tte $\varphi:R-R'$ dans $\Delta$, les foncteurs $u\longmapsto u\circ\varphi$ de $\underline{\mathrm{Hom}}_\Delta(R',E)\longrightarrow\underline{\mathrm{Hom}}_\Delta(R,E)$. (les flèches entre $E$ \unsure{restant|restent} des foncteurs $d$-exacts)

\[
\begin{aligned}
a_0)&\ \text{sont 0-fidèles si }\varphi\text{ est un } f_\lambda^{\unsure{0}}\\
a_1)&\ \deleted{\text{sont}}\ \unsure{\text{1-fid}|\text{pl-fid}|\text{bijectifs}|\text{injectifs}}\ \deleted{\unsure{\text{si }\varphi\text{ est un}}}\quad f_\lambda^1\\
a_2)&\ \text{'' 2-fid}\quad\text{--- --- --}\quad f_\lambda^2
\end{aligned}
\]

\uline{Théorème} \inserted{Avec les notations précédentes,} Soit $(C,\delta)$ comme dans le th. Alors on peut trouver un foncteur \inserted{$\delta$-exact} $C\overset{i}{-}\bar C$, avec $\bar C\in\mathrm{Ob}\,\Delta_f$, et qui soit \inserted{$2$-}universel pour ces propriétés. De plus, $\bar C$ est ess. petite.

[Cela résoud en particulier les Pb 1 et Pb 2 \inserted{et Pb $2'$} pour $\Delta_f$]

% archive sheet number: 6
